\section{Related Work}
\label{sec:related-work}

\subsection{Programmable Filesystem Boundaries}
FUSE is the closest general-purpose mechanism for programmable filesystem
behavior~\cite{linux_fuse}. ExtFUSE moves selected FUSE handling into the
kernel to reduce overhead~\cite{extfuse}. \namei targets a different boundary
because it does not accelerate a filesystem service or implement filesystem
methods. It exposes only state-transition object selection and visibility
before lower filesystem lookup continues.

Wrapfs and related stackable approaches show the long-standing value of adding
filesystem behavior beneath existing applications~\cite{wrapfs}. Bento lowers
the cost and risk of writing in-kernel filesystems~\cite{bento}. These systems
remain relevant baselines and related work, but they ask the developer to own a
broader filesystem interface than \namei's policy boundary.
FUSE and stackable systems own a filesystem service or filesystem methods,
while \namei owns only the object-selection policy decision inside VFS name
resolution.

\subsection{Namespace And Materialization Boundaries}
Bind mounts, projected volumes, symlink forests, and OverlayFS are practical
ways to present alternate views~\cite{kubernetes_projected_volumes,linux_overlayfs}.
They are appropriate when a view can be constructed or layered outside the
transition. \namei instead measures whether object selection can remain a
policy-only boundary under the same oracle.
Namespace mechanisms materialize or layer a view, while \namei updates policy
state and lets each lookup or enumeration event select the lower object for the
current state.

\subsection{Agent Workspace Filesystems}
AgentFS, BranchFS, Sandlock, and YoloFS show that agent workspaces need
branching, copy-on-write state, snapshot-like behavior, and constrained file
views~\cite{agentfs_repo,branchfs_repo,sandlock_repo,yolofs_repo}. These
systems supply transition shapes and oracles. \namei reuses a transition only
when its oracle-relevant effect is selecting or hiding existing lower objects.
Broader sandbox, runtime, snapshot, or custom-filesystem behavior remains part
of the source system's boundary.
The source systems define agent workspace semantics, while \namei tests only
the classified path-view slice of a selected transition.

\subsection{Agent Environment And Build Workloads}
OpenHands, SWE-agent, SWE-ReX, SWE-Factory-Gym, MEnvData-SWE, and SWE-rebench
V2 provide repository repair, Docker-backed environment construction, and
build/test oracles~\cite{openhands_sdk_repo,swe_agent_repo,swe_rex_repo,swe_factory_repo,menvagent_repo,swe_rebench_v2_repo}.
These systems motivate the environment/cache transition in this proposal. The
paper's boundary is narrower because it asks whether the path-view part of a
selected transition can be implemented as policy-only name resolution while the
source system or lower filesystem continues to own environment construction,
command execution, and data-path behavior. These workloads supply environment
and build oracles, while \namei contributes the name-resolution policy boundary
under those oracles.
